\documentclass[
  12pt,
  oneside,
  a4paper,
  chapter=TITLE,
  english,
  portuguese
]{abntex2}

% --- CODIFICAÇÃO ---
\usepackage[utf8]{inputenc}
\usepackage[T1]{fontenc}

% -------------------------------------------------------------------
% FONTE  —  controlada por flags booleanas geradas pelo build.sh
%
%   font-arial     → Liberation Sans     (métricas idênticas à Arial)
%   font-helvetica → Nimbus Sans         (Helvetica livre, URW)
%   font-ubuntu    → Ubuntu              (família clássica, estáticas disponíveis)
%   font-times     → TeX Gyre Termes     (clone livre do Times New Roman)
%
% Padrão (nenhuma flag): TeX Gyre Termes — ABNT NBR 14724 recomendada
% -------------------------------------------------------------------
\usepackage{fontspec}
$if(font-arial)$
  \setmainfont{Liberation Sans}
$else$$if(font-helvetica)$
  \setmainfont{Nimbus Sans}
$else$$if(font-ubuntu)$
  \setmainfont{Ubuntu}
$else$
  \setmainfont{TeX Gyre Termes}
$endif$$endif$$endif$

% --- IDIOMA (XeLaTeX-nativo, integra com fontspec) ---
\usepackage{polyglossia}
\setdefaultlanguage[variant=brazilian]{portuguese}

% --- MARGENS (ABNT NBR 10719) ---
\usepackage[a4paper,top=3cm,left=3cm,right=2cm,bottom=2cm]{geometry}

% --- ESPAÇAMENTO (1,5 entrelinhas) ---
\usepackage{setspace}
\OnehalfSpacing

% --- PACOTES ESSENCIAIS ---
\usepackage{graphicx}
\usepackage{indentfirst}
\usepackage{longtable}
\usepackage{booktabs}
\usepackage{array}
\usepackage{calc}
\usepackage{amsmath,amssymb,amsfonts}
\usepackage{textcomp}
\usepackage{color,xcolor}
\usepackage{listings}
\usepackage{fancyvrb}
\usepackage{upquote}
\usepackage{csquotes}
\usepackage{microtype}
\usepackage{hyperref}
\hypersetup{
  colorlinks=true,
  linkcolor=black,
  filecolor=black,
  urlcolor=black,
  citecolor=black,
}
% --- FIGURE PLACEMENT ---
% LaTeX's default float algorithm ([htbp]) repositions figures freely,
% causing images to appear far from their insertion point and overlap
% paragraphs. [H] forces each figure to appear exactly where declared.
\usepackage{float}
\usepackage{caption}
\floatplacement{figure}{H}

% --- SUPORTE A CÓDIGO (Pandoc) ---
\DefineVerbatimEnvironment{Highlighting}{Verbatim}{commandchars=\\\{\}}
\newenvironment{Shaded}{}{}
\newcommand{\passthrough}[1]{#1}
\newcommand{\KeywordTok}[1]{\textcolor{blue}{\textbf{#1}}}
\newcommand{\ControlFlowTok}[1]{\textcolor{blue}{\textbf{#1}}}
\newcommand{\StringTok}[1]{\textcolor{blue}{#1}}
\newcommand{\CommentTok}[1]{\textcolor{gray}{\textit{#1}}}
\newcommand{\FunctionTok}[1]{\textcolor{cyan}{#1}}
\newcommand{\NormalTok}[1]{#1}
\newcommand{\OtherTok}[1]{#1}
\newcommand{\ErrorTok}[1]{\textcolor{red}{#1}}
\newcommand{\DataTypeTok}[1]{\textcolor{purple}{#1}}
\newcommand{\DecValTok}[1]{\textcolor{orange}{#1}}
\newcommand{\AttributeTok}[1]{\textcolor{orange}{#1}}
\newcommand{\BuiltInTok}[1]{\textcolor{purple}{#1}}
\newcommand{\ImportTok}[1]{\textcolor{teal}{#1}}
\newcommand{\VariableTok}[1]{\textcolor{black}{#1}}
\newcommand{\OperatorTok}[1]{#1}
\newcommand{\ExtensionTok}[1]{#1}
\newcommand{\SpecialCharTok}[1]{\textcolor{violet}{#1}}
\newcommand{\DocumentationTok}[1]{\textcolor{gray}{\itshape #1}}
\newcommand{\AnnotationTok}[1]{\textcolor{red}{\textbf{#1}}}
\newcommand{\CommentVarTok}[1]{\textcolor{gray}{\itshape #1}}
\newcommand{\ConstantTok}[1]{\textcolor{red}{#1}}
\newcommand{\SpecialStringTok}[1]{\textcolor{violet}{#1}}
\newcommand{\VerbatimStringTok}[1]{\textcolor{green}{#1}}
\newcommand{\CharTok}[1]{\textcolor{green}{#1}}
\newcommand{\BaseNTok}[1]{\textcolor{orange}{#1}}
\newcommand{\FloatTok}[1]{\textcolor{orange}{#1}}
\newcommand{\PreprocessorTok}[1]{\textcolor{purple}{\textbf{#1}}}
\newcommand{\RegionMarkerTok}[1]{#1}
\newcommand{\InformationTok}[1]{\textcolor{gray}{\itshape #1}}
\newcommand{\WarningTok}[1]{\textcolor{red}{\textbf{#1}}}
\newcommand{\AlertTok}[1]{\textcolor{red}{\textbf{#1}}}

% -------------------------------------------------------------------
% METADADOS
% -------------------------------------------------------------------
\titulo{$title$}
\autor{$for(author)$$if(author.name)$$author.name$$else$$author$$endif$$sep$\\$endfor$}
\instituicao{$if(institution)$$institution$$else$Instituição de Ensino$endif$}
\local{$if(city)$$city$$else$Cidade$endif$}
\data{$if(date)$$date$$else$\today$endif$}

% -------------------------------------------------------------------
% NORMALIZAÇÃO DOS TÍTULOS
% ABNT NBR 14724 specifies 12 pt for all section headings.
% abntex2 defaults to sizes larger than \normalsize for \section and
% below; override all levels to match the chapter heading size.
% -------------------------------------------------------------------
\renewcommand{\ABNTEXchapterfont}{\bfseries}
\renewcommand{\ABNTEXchapterfontsize}{\normalsize}
\renewcommand{\ABNTEXsectionfontsize}{\normalsize}
\renewcommand{\ABNTEXsubsectionfontsize}{\normalsize}
\renewcommand{\ABNTEXsubsubsectionfontsize}{\normalsize}
\renewcommand{\ABNTEXsubsubsubsectionfontsize}{\normalsize}

% -------------------------------------------------------------------
\begin{document}
\frenchspacing

% =====================================================================
% CAPA
% =====================================================================
\begin{capa}
  \center

  $if(logo)$
  \includegraphics[width=$if(logo-width)$$logo-width$$else$0.12$endif$\textwidth]{$logo$}\\[0.3cm]
  $endif$

  {\ABNTEXchapterfont\bfseries\large
    \MakeUppercase{$if(institution)$$institution$$else$INSTITUIÇÃO DE ENSINO$endif$}\\
    $if(department)$\MakeUppercase{$department$}\\$endif$
    $if(course)$\MakeUppercase{$course$}$endif$
  }\\[2.5cm]

  {\ABNTEXchapterfont\bfseries\large
    $for(author)$$if(author.name)$\MakeUppercase{$author.name$}$else$\MakeUppercase{$author$}$endif$$sep$\\$endfor$
  }\\[4cm]

  {\ABNTEXchapterfont\bfseries\large
    \MakeUppercase{$title$}$if(subtitle)$: $subtitle$$endif$
  }\\[1cm]

  $if(report-type)$
  {\large $report-type$}\\[\fill]
  $else$
  \vfill
  $endif$

  {\ABNTEXchapterfont\large
    \MakeUppercase{$if(city)$$city$$else$CIDADE$endif$}$if(state)$ -- $state$$endif$\\
    $if(date)$$date$$else$\today$endif$
  }
\end{capa}

% =====================================================================
% FOLHA DE ROSTO
% =====================================================================
\begin{folhaderosto}
  \center

  {\ABNTEXchapterfont\bfseries\large
    $for(author)$$if(author.name)$\MakeUppercase{$author.name$}$else$\MakeUppercase{$author$}$endif$$sep$\\$endfor$
  }\\[15\baselineskip]

  {\ABNTEXchapterfont\bfseries\large
    \MakeUppercase{$title$}$if(subtitle)$: $subtitle$$endif$
  }\\[2cm]

  \hspace*{8cm}
  \parbox{8cm}{\normalsize\setlength{\parskip}{0pt}\setlength{\parindent}{0pt}\normalfont
    $if(nature)$$nature$$else$%
      $if(report-type)$$report-type$.$else$Relatório técnico.$endif$
      $if(course)$ Curso de $course$.$endif$
      $if(institution)$ $institution$$if(campus)$, $campus$$endif$.$endif$
    $endif$
  }\\[0.8cm]

  $if(supervisor)$
  \hspace*{8cm}
  \parbox{8cm}{\normalsize\normalfont
    $if(supervisor-label)$$supervisor-label$$else$Supervisor(a)$endif$: $supervisor$
  }\\[0.5cm]
  $endif$

  \vfill

  {\ABNTEXchapterfont\large
    \MakeUppercase{$if(city)$$city$$else$CIDADE$endif$}$if(state)$ -- $state$$endif$\\
    $if(date)$$date$$else$\today$endif$
  }
\end{folhaderosto}

$if(abstract-pt)$
\newpage
\thispagestyle{empty}
\begin{resumoumacoluna}
\noindent\textbf{RESUMO:}\vspace{\onelineskip}

$abstract-pt$

\vspace{\onelineskip}
\noindent\textbf{Palavras-chave}: $if(keywords-pt)$$keywords-pt$$endif$.
\end{resumoumacoluna}
\clearpage
$endif$

$if(list-of-figures)$
\pdfbookmark[0]{\listfigurename}{lof}
\listoffigures*
\cleardoublepage
$endif$

$if(list-of-tables)$
\pdfbookmark[0]{\listtablename}{lot}
\listoftables*
\cleardoublepage
$endif$

\pdfbookmark[0]{\contentsname}{toc}
\tableofcontents*
\cleardoublepage

\textual
\setcounter{secnumdepth}{3}
% Fix: \textual resets internal spacing state in abntex2, causing the first
% paragraph of the body to lose its indentation. Restoring \parindent here
% ensures consistent indentation across all body paragraphs.
\setlength{\parindent}{1.3cm}

$body$

\postextual

$if(bibliography)$
\bibliographystyle{abntex2-alf}
\bibliography{$bibliography$}
$endif$

$if(appendix)$
\begin{apendicesenv}
$appendix$
\end{apendicesenv}
$endif$

\end{document}
